\documentclass[12pt,a4paper]{article}

% ------------------------------------------------
% Pacotes
% ------------------------------------------------
\usepackage[utf8]{inputenc}
\usepackage[T1]{fontenc}
\usepackage[brazil]{babel}
\usepackage{graphicx}
\usepackage[table]{xcolor}
\usepackage{array}
\usepackage{tabularx}
\usepackage{geometry}
\usepackage{fancyhdr}
\usepackage{mathptmx} % Times New Roman
\usepackage{enumitem}
\usepackage{amsfonts}
\usepackage{amsmath} % Recomendado para o comando \text{}
\usepackage{amssymb}  
% ------------------------------------------------ 
% Configuração de página
% ------------------------------------------------
\geometry{
    a4paper,
    left=2cm,
    right=2cm,
    top=2cm,
    bottom=2cm,
    headheight=15pt
}
% Caixa para resposta discursiva
\newcommand{\resposta}[1]{
\noindent
\framebox[\textwidth]{
\begin{minipage}{0.98\textwidth}
\vspace{#1}
\end{minipage}
}
}

\definecolor{forestgreen}{RGB}{0,100,50}
\pagestyle{fancy}
\fancyhf{} % limpa cabeçalho e rodapé

% Linha do cabeçalho e rodapé
\renewcommand{\headrulewidth}{2pt}
\renewcommand{\footrulewidth}{2pt}

\renewcommand{\headrule}{\hbox to\headwidth{%
  \color{forestgreen}\leaders\hrule height \headrulewidth\hfill}}

\renewcommand{\footrule}{\hbox to\headwidth{%
  \color{forestgreen}\leaders\hrule height \footrulewidth\hfill}}
% Ajuste de altura do cabeçalho
\setlength{\headheight}{33pt}

% Cabeçalho com imagem
\fancyhead[L]{%
  \includegraphics[height=1.2cm]{Universidade_de_Rio_Verde.jpg}
}

\fancyhead[R]{%
  \raggedleft
  \textcolor{forestgreen}{%
    \small
    \textbf{Lista de Exercícios — Provas Matemáticas}\\
    Me.~Sergio Souza Novak\\
    \textit{Lógica para Engenharia de Software}
  }
}

% Rodapé
\fancyfoot[C]{%
  \textcolor{forestgreen}{\thepage}
}

\setlength{\parindent}{0pt}

% ------------------------------------------------
% Comandos auxiliares
% ------------------------------------------------
\newcolumntype{L}[1]{>{\raggedright\arraybackslash}p{#1}}
\newcolumntype{C}[1]{>{\centering\arraybackslash}p{#1}}

% ------------------------------------------------
% Documento
% ------------------------------------------------
\begin{document}

% -------------------------------------


\textbf{1)} Construa a tabela-verdade completa da proposição $P \Leftrightarrow Q$. Em seguida, indique explicitamente em quais casos a dupla implicação é verdadeira.

\vspace{0.2cm}
\resposta{8cm}

\vspace{0.8cm}

\textbf{2)} Considere a implicação: \textit{“Se um triângulo é equilátero, então ele é isósceles.”} Identifique claramente a hipótese e a conclusão.

\vspace{0.2cm}
\resposta{3cm}

\vspace{0.8cm}

\textbf{3)} Demonstre que $A \cup A = A$.

\vspace{0.2cm}
\resposta{5cm}
 
\vspace{0.8cm}

\textbf{4)} Seja $m$ um inteiro par e $n$ um inteiro par. Prove que a soma $m+n$ é um número par.

\vspace{0.2cm}
\resposta{4cm}

\vspace{0.8cm}

\textbf{5) } Uma relação de equivalência é uma relação binária $R$ em um conjunto $A$, tal que $R$ é reflexiva, simétrica e transitiva.

Considere as relações binárias apresentadas a seguir:
\begin{itemize}
    \item $R_1 = \{(a, b) : a, b \in \mathbb{Z} \text{ e } a = b\}$;
    \item $R_2 = \{(a, b) : a, b \in \mathbb{Z} \text{ e } a < b\}$;
    \item $R_3 = \{(a, b) : a, b \in \mathbb{Z} \text{ e } a = b + 1\}$;
    \item $R_4 = \{(a, b) : a, b \in \mathbb{Z} \text{ e } a - b \text{ é um múltiplo de 3}\}$.
\end{itemize}

São relações de equivalência apenas o que se apresenta em:
\begin{enumerate}[label=\Alph*)]
    \item $R_1$ e $R_4$.
    \item $R_2$ e $R_3$.
    \item $R_1$ e $R_3$. 
    \item $R_1, R_2$ e $R_4$.
    \item $R_2, R_3$ e $R_4$.
\end{enumerate} 

\vspace{0.8cm}

\textbf{6)} Prove que $A \cap A = A$.

\vspace{0.2cm}
\resposta{5cm}

\vspace{0.8cm}

\textbf{7)} Sejam $x$ e $y$ dois inteiros pares. Prove que o produto $xy$ é um múltiplo de 4.

\vspace{0.2cm}
\resposta{4cm}

\vspace{0.8cm}

\textbf{8) } Uma relação $R$ em um conjunto $S$ é uma relação de equivalência se ela satisfizer todas as propriedades a seguir:
\begin{enumerate}
    \item \textbf{Reflexividade.} Para todo $a \in S$, $a R a$;
    \item \textbf{Simetria.} Para todos $a, b \in S$, $a R b \Leftrightarrow b R a$;
    \item \textbf{Transitividade.} Para todos $a, b, c \in S$, se $a R b$ e $b R c$, então, $a R c$.
\end{enumerate}

Uma partição de um conjunto $S$ é uma coleção de subconjuntos disjuntos não vazios, cuja união é igual a $S$. Se $R$ é uma relação de equivalência em um conjunto $S$ e se $x \in S$, denota-se por $[x]$ o conjunto de todos os elementos relacionados a $x$ em $S$ e chama-se esse conjunto de classe de equivalência de $x$.

A partir dessas informações, considere que, em uma rede social corporativa, haja 8 usuários de diferentes departamentos. Os administradores dessa rede pretendem mapear os usuários que compartilham os mesmos interesses profissionais. Assim, fizeram uma tabela relacionando os usuários considerando pares. Dessa forma, dois usuários fazem parte de uma linha na tabela se possuem o mesmo interesse profissional. Obviamente, apesar de não estar evidente na tabela, um usuário também tem relação consigo mesmo.

Os usuários estão numerados de 1 até 8 e a tabela resultante pode ser vista a seguir.

\begin{center}
\textbf{Usuários com interesses semelhantes} \\
\vspace{0.2cm}
\begin{tabular}{|c|c|}
\hline
1 & 3 \\ \hline
2 & 5 \\ \hline
3 & 8 \\ \hline
5 & 7 \\ \hline
2 & 6 \\ \hline
4 & 4 \\ \hline
8 & 1 \\ \hline
6 & 5 \\ \hline
7 & 2 \\ \hline
1 & 8 \\ \hline
\end{tabular}
\end{center}

Com base nesse cenário e no conceito de relações de equivalência, assinale a opção correta.
 
\begin{enumerate}[label=\Alph*)]
    \item A relação descrita pela tabela não é uma relação de equivalência.
    \item O subconjunto de usuários representados pelos números 2, 5, 6 e 7 forma uma classe de equivalência.
    \item A relação descrita pela tabela apresenta a propriedade de simetria, mas não a de transitividade.
    \item O subconjunto de usuários representados pelos números 1, 3 e 4 forma uma classe de equivalência.
    \item O usuário 4 não pertence a nenhuma classe de equivalência.
\end{enumerate}

% ------------------------------------------------
% Colinha / Resumo de Definições
% ------------------------------------------------

\section*{Resumo de Definições (Referência)}

\textbf{Conjunto Vazio}

\[
\varnothing = \{x \mid x \neq x\}
\]

\[
(\forall x)(x \notin \varnothing)
\]

\textbf{Paridade}

\[
n \in \mathbb{Z} \text{ é par } \Leftrightarrow \exists k \in \mathbb{Z} \; (n = 2k)
\]

\[
n \in \mathbb{Z} \text{ é ímpar } \Leftrightarrow \exists k \in \mathbb{Z} \; (n = 2k + 1)
\]

\textbf{Múltiplo de 4}

\[
n \in \mathbb{Z} \text{ é múltiplo de 4 } \Leftrightarrow \exists k \in \mathbb{Z} \; (n = 4k)
\]

\vspace{0.3cm}

\textbf{Relação de Contido}

\[
A \subseteq B \Leftrightarrow (\forall x)(x \in A \Rightarrow x \in B)
\]

\vspace{0.3cm}

\textbf{União de Conjuntos}

\[
A \cup B = \{x \mid x \in A \lor x \in B\}
\]

\vspace{0.3cm}

\textbf{Interseção de Conjuntos}

\[
A \cap B = \{x \mid x \in A \land x \in B\}
\]

\vspace{0.3cm}

\textbf{Relação Binária}

Uma relação $R$ entre dois conjuntos $A$ e $B$ é um subconjunto do produto cartesiano:

\[
R \subseteq A \times B
\]

Quando $A=B$, dizemos que $R$ é uma relação em $A$:

\[
R \subseteq A \times A
\]

\vspace{0.3cm}

\textbf{Propriedades de Relações}

\textbf{Reflexiva}

\[
(\forall a \in A)(aRa)
\]

\textbf{Simétrica}

\[
(\forall a,b \in A)(aRb \Rightarrow bRa)
\]

\textbf{Transitiva}

\[
(\forall a,b,c \in A)((aRb \land bRc) \Rightarrow aRc)
\]
\end{document}
